\documentclass[pdf, aspectratio=169, 12pt]{beamer}
\usepackage[]{hyperref, graphicx, siunitx, lmodern, tikz, booktabs, physics, multicol}
\usepackage[mode=buildnew]{standalone}
\usepackage{pdfpc-commands}
\usepackage{pgfplots}
\pgfplotsset{compat=1.16}

\usetheme{Python}

\graphicspath{ {Images/} }

\sisetup{per-mode=symbol}
\usetikzlibrary{calc, patterns, decorations.markings, decorations.pathmorphing, shapes, fit}

%Preamble
\title{Not for the individual, but for the class...}
\author{Jed Rembold}
\date{April 8, 2020}

\begin{document}

\begin{frame}{Announcements}
	\begin{itemize}
		\item Homework
			\begin{itemize}
				\item Homework 9 is out!
				\item Homework 8 is due on Friday as well!
			\end{itemize}
		\item Virtual CS Tea tomorrow!
			\begin{itemize}
				\item Connection info posted on Campuswire
				\item Just socializing and catching up this week
			\end{itemize}
		\item CS/Data Science professor candidate talks upcoming
			\begin{itemize}
				\item Teaching talk today 3-4
				\item Research talk tomorrow 4:10-5:30
				\item Zoom info on Campuswire
			\end{itemize}
			
		\item Polling: \nolinkurl{rembold-class.ddns.net}
	\end{itemize}
\end{frame}

\begin{frame}[fragile]{Review Question}
	\begin{columns}
		\column{0.5\textwidth}
		What would be the output of \pyi{print(D.x)} for the very contrived bit of code to the right?
		\begin{poll}
		\item 15
		\item 13
		\item 10
		\item \pyi{None}
		\end{poll}
		\exsol{15}
		
		\column{0.5\textwidth}
		\scriptsize
		\vspace{-8mm}
		\begin{pythoncode}
			class ObjA:
				def __init__(self):
					self.x = 5

			class ObjB(ObjA):
				def __init__(self):
					ObjA.__init__(self)
					self.x = 8

			class ObjC(ObjB):
				def __init__(self):
					ObjA.__init__(self)
					self.y = self.x + 2

			class ObjD(ObjC):
				def __init__(self):
					ObjC.__init__(self)
					ObjB.__init__(self)
					self.x += self.y

			D = ObjD()
		\end{pythoncode}
	\end{columns}
\end{frame}

\begin{frame}{Multiple Inheritance}
	\begin{itemize}
		\item A class is not limited to inheriting from only a single parent!
		\item Can provide multiple parents!
			\begin{itemize}
				\item Ordering matters!
			\end{itemize}
		\item Example: a class that forms the intersection of two other types
	\end{itemize}
\end{frame}

\begin{frame}[fragile]{Pieces of Inheritance}
	\begin{columns}
		\column{0.5\textwidth}
	\begin{center}
		\begin{tikzpicture}[
			class/.style={draw, very thick, rounded corners, minimum size=2cm},
			att/.style={font=\scriptsize, text width=2cm, align=center},
			methods/.style={inner sep=0pt, outer sep=0pt, minimum width=1.5cm, minimum height=3mm},
			]

			\coordinate (A) at (0,0);
			\node[att, Blue, above = 2mm of A.center](At) {attributes};
			\node[methods, fill=Blue, below = 2mm of A.center](Ab) {};
			\node[class,Blue, fit=(At)(Ab)](Ac) {};

			\coordinate[right = 4cm of A] (B);
			\node[att, Green, above = 5mm of B.center](Bt) {attributes};
			\node[att, Blue, above = 2mm of B.center] {attributes};
			\node[methods, fill=Green, below = 2mm of B.center] {};
			\node[methods, fill=Blue, below = 6mm of B.center](Bb) {};
			\node[class,Green, fit={(Bt)(Bb)}](Bc) {};

			\coordinate[below = 3.5cm of A] (C);
			\node[att, Red, above = 2mm of C.center](Ct) {attributes};
			\node[methods, fill=Red, below = 2mm of C.center](Cb) {};
			\node[class,Red, fit=(Ct)(Cb)](Cc) {};

			\coordinate[right = 4cm of C] (D);
			\node[att, Orange, above = 11mm of D.center](Dt) {attributes};
			\node[att, Red, above = 8mm of D.center] {attributes};
			\node[att, Green, above = 5mm of D.center] {attributes};
			\node[att, Blue, above = 2mm of D.center] {attributes};
			\node[methods, fill=Orange, below = 2mm of D.center] {};
			\node[methods, fill=Red, below = 6mm of D.center] {};
			\node[methods, fill=Green, below = 10mm of D.center] {};
			\node[methods, fill=Blue, below = 14mm of D.center](Db) {};
			\node[class,Orange, fit={(Dt)(Db)}](Dc) {};

			\path[ultra thick, -stealth, Purple]
				(Dc) edge (Cc) edge (Bc)
				(Bc) edge (Ac);

			\only<2->{
				\node[methods, fill=Green, below = 6mm of D.center] {};
				\node[methods, fill=Blue, below = 10mm of D.center] {};
				\node[methods, fill=Red, below = 14mm of D.center] {};
			}
			\only<3>{
				\node[att, Green, above = 8mm of D.center] {attributes};
				\node[att, Blue, above = 5mm of D.center] {attributes};
				\node[att, Red, above = 2mm of D.center] {attributes};
			}
			\only<4>{
				\node[att, Green, above = 11mm of D.center] {attributes};
				\node[att, Blue, above = 8mm of D.center] {attributes};
				\node[att, Orange, above = 5mm of D.center] {attributes};
				\node[att, Red, above = 2mm of D.center] {attributes};
			}

		\end{tikzpicture}
	\end{center}
		
		\column{0.5\textwidth}
		\begin{onlyenv}<1>
			\begin{pythoncode}
				class Orange(Red, Green):
					def __init__(self):
						Green.__init__.self()
						Red.__init__.self()
						self.attributes
			\end{pythoncode}
		\end{onlyenv}
		\begin{onlyenv}<2>
			\begin{pythoncode}
				class Orange(Green, Red):
					def __init__(self):
						Green.__init__.self()
						Red.__init__.self()
						self.attributes
			\end{pythoncode}
		\end{onlyenv}
		\begin{onlyenv}<3>
			\begin{pythoncode}
				class Orange(Green, Red):
					def __init__(self):
						Red.__init__.self()
						Green.__init__.self()
						self.attributes
			\end{pythoncode}
		\end{onlyenv}
		\begin{onlyenv}<4>
			\begin{pythoncode}
				class Orange(Green, Red):
					def __init__(self):
						Red.__init__.self()
						self.attributes
						Green.__init__.self()
			\end{pythoncode}
		\end{onlyenv}
	\end{columns}
\end{frame}

\begin{frame}{Thought Experiment}
	\begin{columns}
		\column{0.5\textwidth}
		Suppose you wanted to program the background logic for playing a game of checkers. What different classes might you want to implement? What about any sub/super classes?
		
		\vspace{5mm}
		List them and their relationships (if any) out. We'll talk about them in a moment.
		
		\column{0.5\textwidth}
		\begin{center}
			\begin{tikzpicture}[
				block/.style={minimum size=1cm, draw, white},
				scale=0.75,
				transform shape,
				]
				\foreach \x in {1,3,...,8}{
					\foreach \y in {1,3,...,8}{
						\node[block, fill=Red] at (\x,\y) {};
						\node[block, fill=BG] at (\x, \y+1) {};
						\node[block, fill=Red] at (\x+1,\y+1) {};
						\node[block, fill=BG] at (\x+1, \y) {};
					}
				}
			\end{tikzpicture}
		\end{center}
	\end{columns}
\end{frame}

\begin{frame}{Variable Types}
	\begin{itemize}
		\item So far we have been looking at \alert{instance variables}.
			\begin{itemize}
				\item A data attribute that is associated to a particular instance or object
			\end{itemize}
		\item Much of the time, this is what we want when dealing with OOP ideas
		\item Sometimes though it is useful to create an attribute that is tied directly to the class itself, and not to a specific object.
			\begin{itemize}
				\item Called a \alert{class variable}.
			\end{itemize}
	\end{itemize}
\end{frame}

\begin{frame}[fragile]{Creating a class variable}
	\begin{itemize}
		\item No usage of \pyi{self}, this variable applies to all!
		\item Define just like you would a normal variable in a function
		\item Usually defined at the top before the methods
	\end{itemize}
	
	\begin{pythoncode}
		class MyBestClass:
			a_class_variable = 'hello!'

			def __init__(self):
				self.a_instance_variable = 'Hi!'
	\end{pythoncode}
\end{frame}

\begin{frame}[fragile]{Accessing a class variable}
	\vspace{5mm}
	\begin{itemize}
		\item Can access directly by referencing the class:
			\begin{itemize}
				\item \pyi{print(MyBestClass.a_class_variable)}
				\item Requires no instance to actually exist
			\end{itemize}
		\item Can access from an instance just like any other instance attribute
			\begin{pythoncode}
				A = MyBestClass()
				print(A.a_class_variable)
			\end{pythoncode}
		\item If changing or setting a class variable, \alert{must} use the \pyi{ClassName.Class_Variable} notation!
			\begin{itemize}
				\item Using \pyi{self} or the object will create an instance variable instead!
			\end{itemize}
	\end{itemize}
	\pause
	\begin{columns}
		\column{0.5\textwidth}
		\begin{pythoncode}
			MyBestClass.class_var = 4
		\end{pythoncode}
		\begin{tikzpicture}[overlay, remember picture]
			\draw[ultra thick, Green, rounded corners] (-.2,0.3) rectangle + (7,1);
		\end{tikzpicture}
		
		\column{0.4\textwidth}
		\begin{pythoncode}
			A.class_var = 4
		\end{pythoncode}
		\begin{tikzpicture}[overlay, remember picture]
			\draw[ultra thick, Red] (-.2,0.3) -- + (4.5,1) (-.2,1.3) -- +(4.5,-1);
		\end{tikzpicture}
	\end{columns}
\end{frame}

\begin{frame}{When to Use}
	\begin{itemize}
		\item Do all members or objects of your class share a common trait?
			\begin{itemize}
				\item Think about if it would make sense for one member of that class to have that trait altered.
				\item If no, might be a good candidate for a class variable
			\end{itemize}
		\item Are there specific constants or known values that all instances of your class will rely on?
		\item Do you need some sort of global counter to keep track of the amount of instances you have created?
	\end{itemize}
	
\end{frame}




\end{document}

